\section{Parameter Estimation}

%------------------------------------------------

\subsection{Physical Parameters}

\begin{frame}
    \frametitle{Assumptions for Physical Parameters Estimation}

    Assumptions
    \begin{itemize}
        \item the pendulum is always the same
        \item physical parameters don't change over time
    \end{itemize}

    Known
    \begin{itemize}
        \item dynamics of the system 
        \item state at each time step
        \item some physical parameters (optional)
    
    Goal
    \begin{itemize}
        \item estimate unknown physical parameters
    \end{itemize}

\end{frame}

\begin{frame}
    \frametitle{Architecture for Physical Parameters Estimation}

    We propose to use an integrator plus a dynamic model with learnable physical parameters:

    \begin{figure}
        \centering
        \includegraphics[width=0.4\textwidth]{images/integrator.png}
    \end{figure}

    Training recipe:
    \begin{itemize}
        \item loss (MSE, MAE)
        \item curricular steps (increasing h)
        \item data transformation (sin-cos)
    \end{itemize}
    
\end{frame}

\begin{frame}
    \frametitle{Architecture for Physical Parameters Estimation}

    \begin{figure}
        \centering
        \includegraphics[width=0.4\textwidth]{images/dyn_ode.png}
    \end{figure}

    Fixing some parameters can help the estimation of the others.

    \begin{figure}
        \centering
        \includegraphics[width=0.4\textwidth]{images/dyn_ode_fixed.png}
    \end{figure}
    
\end{frame}

\begin{frame}
    \frametitle{Results for Physical Parameters Estimation}

    Fixing parameters make the difference.

    Gauge invariance issue can affect the estimation.

    \begin{figure}
        \centering
        \includegraphics[width=0.7\textwidth]{images/phys_params_results.png}
    \end{figure}
\end{frame}

%------------------------------------------------

\subsection{Dynamics Matrices Entries}

\begin{frame}
    \frametitle{Assumptions for Matrices Entries Estimation}

    Assumptions
    \begin{itemize}
        \item the pendulum is always the same
        \item physical parameters don't change over time
    \end{itemize}

    Known
    \begin{itemize}
        \item dynamics of the system
        \item state at each time step 
        \item some matrices entries (optional)
    \end{itemize}
    
    Goal
    \begin{itemize}
        \item estimate unknown matrices entries
    \end{itemize}

\end{frame}

\begin{frame}
    \frametitle{NN for Matrices Entries Estimation}

    The first approach we propose is to use a NN that simply learns the entries of the matrices.

    Requirements
    \begin{itemize}
        \item GT values for matrices in dataset
    \end{itemize}

    \begin{figure}
        \centering
        \includegraphics[width=0.4\textwidth]{images/entries_NN.png}
    \end{figure}


\end{frame}

\begin{frame}
    \frametitle{UDE for Matrices Entries Estimation}

    The second approach is a UDE dynamic layer combinaed with the integrator

    \begin{figure}
        \centering
        \includegraphics[width=0.4\textwidth]{images/dyn_UDE.png}
    \end{figure}

\end{frame}

\begin{frame}
    \frametitle{Informed UDE for Matrices Entries Estimation}

    We can enhance the UDE approach `informing' better our UDE

    \begin{figure}
        \centering
        \includegraphics[width=0.4\textwidth]{images/informed_UDE.png}
    \end{figure}

     \begin{figure}
        \centering
        \includegraphics[width=0.4\textwidth]{images/better_informed_UDE.png}
    \end{figure}

\end{frame}

%------------------------------------------------

\subsection{Results}


\begin{frame}
    \frametitle{Results for Physical Parameters Estimation}

    Learning directly on the entries makes the task easy.

    Gauge invariance issue can affect the estimation using UDE. 

    Informing better UDE helps a lot. 

    \begin{figure}
        \centering
        \includegraphics[width=0.7\textwidth]{images/matirces_entries_results.png}
    \end{figure}
\end{frame}



